\section{TabPFN 与 ExtraTrees 融合方法}
\subsection{条件回归与原始模型}
记第 $i$ 条记录的药物和溶剂结构为 $d_i$ 与 $s_i$，温度为 $T_i$，母体分组为 $g_i$。模型使用两种结构和 CSV 原始温度预测对数溶解度 $y_i$。同一药物--溶剂组合在多个温度下有观测时，两种分子结构保持固定。原始模型 SOLAR-H 利用这一点，将预测写为
\begin{equation}
\hat y_{{\rm H},i}=b(z_i)+\gamma\,\tau_i\,\beta^\top u_i,\qquad u_i=[1;z_i],
\end{equation}
其中 $z_i$ 为结构特征，$\tau_i$ 为按训练折标准化的温度，$b$ 为 ExtraTrees 结构基线，$\gamma$ 为内层验证选择的温度增益。同一组合 $c$ 的 $z_i$ 固定，在组合内中心化可以消去 $b$。代码据此用有效温度偏差 $|\tau_i-\bar\tau_c|>10^{-9}$ 的记录拟合岭回归，得到 $\beta$：
\begin{equation}
\hat\beta=\arg\min_\beta
\sum_i w_i\big[(y_i-\bar y_{c(i)})-(\tau_i-\bar\tau_{c(i)})\,\beta^\top u_i\big]^2
+\lambda\|\beta\|_2^2 .
\end{equation}
求和仅覆盖上述有效记录，$\bar y_c$ 与 $\bar\tau_c$ 为该组合在训练折内的均值。给定 $\gamma$ 后，树在温度修正后的目标 $y_i-\gamma\tau_i\hat\beta^\top u_i$ 上拟合结构基线。后续的融合方法将温度与分子描述符一并输入两个学习器，由树切分和表格上下文建模条件关系。

\subsection{分子描述符与交互特征}
令 $\phi_k(m)$ 为分子 $m$ 的 $k$ 项描述符。每条记录的基础特征为 $[\phi_{11}(d_i);\phi_{11}(s_i);T_i]$，共 $2\times11+1=23$ 维。额外纳入完整 RDKit 描述符时，令 $D_i=\phi_{217}(d_i)$、$S_i=\phi_{217}(s_i)$，则基础拼接扩展至 $2\times217+1=435$ 维。为显式表示同类描述符的共同变化及差异，在此基础上拼接逐元素乘积与绝对差：
\begin{equation}
x_i^{\rm int}
=[D_i;S_i;T_i;D_i\odot S_i;|D_i-S_i|]
\in\mathbb R^{869}.
\end{equation}
例如，药物和溶剂各自的 MolLogP 除原始值外，还生成乘积和绝对差两列。两种交互各增加 217 维。再加入标准化药物母体结构的 217 项描述符，得到 1086 维的候选表示。树分支对当前训练折拟合缺失描述符的中位数插补和零方差筛选，并将同一变换应用于验证折。

\subsection{按母体加权的 ExtraTrees 候选}
单纯按记录计数时，训练折中含 $n_g$ 条记录的母体会贡献 $n_g$ 次损失。对训练折 $\mathcal T$，记 $n_g(\mathcal T)$ 为母体 $g$ 的记录数，$N_{\mathcal T}$ 和 $G_{\mathcal T}$ 分别为总记录数和母体数。将样本权重设为
\begin{equation}
w_i(q)=\frac{n_{g_i}(\mathcal T)^{-q}}
{N_{\mathcal T}^{-1}\sum_{j\in\mathcal T}n_{g_j}(\mathcal T)^{-q}},
\qquad q\in\{0,1\}.
\end{equation}
当 $q=0$ 时每条记录权重为 1。当 $q=1$ 时，分母为 $G_{\mathcal T}/N_{\mathcal T}$，同一母体的权重之和为 $N_{\mathcal T}/G_{\mathcal T}$。在完整开发集上，每个母体的总权重为 $1507/128\approx 11.77$。

对树节点中的记录集合 $L$，加权平方损失为 $J_L(v)=\sum_{i\in L}w_i(y_i-v)^2$。令 $J_L'(v)=0$，得到叶预测 $\hat v_L=\sum_{i\in L}w_i y_i/\sum_{i\in L}w_i$。对候选切分 $L=L_{\rm left}\cup L_{\rm right}$，比较母节点与两子节点各自最小平方损失的差，选择降幅最大的候选切分。最终输出 300 棵树的叶预测均值。每棵树不限制深度，叶节点至少含 2 条记录，每次切分考虑 0.7 比例的特征。完整、交互及含母体描述符三种表示与两个 $q$ 值组合成六项候选，内层按母体等权绝对误差选择。该指标与树拟合时的平方损失分别用于模型选择和节点拟合。

\subsection{TabPFN 的上下文预测}
TabPFN 分支使用公开 v2 回归权重 $\theta$。对一个验证样本 $x$，其上下文 $\mathcal C_{\mathcal T}=\{(x_j,y_j):j\in\mathcal T\}$ 由当前训练折的描述符、温度及对数溶解度构成。点预测取条件预测分布的均值\cite{tabpfn2025}：
\begin{equation}
f_{\rm PFN}(x)=\mathbb E_{\theta}[Y\mid x,\mathcal C_{\mathcal T}].
\end{equation}
模型权重固定，上下文的训练记录等权输入。内层比较 23 维与 435 维表示，并为每种表示比较 2 或 4 个内部集成成员。

\subsection{内层折外选参和融合}
对每个外层训练集，按母体重新划分三折。每个分支的每项候选均对三折中的验证记录生成折外预测。合并三折后，分别以母体等权 MAE（式\ref{eq:parentmae}）选择两分支配置。记内层记录 $i$ 对应的所选分支预测为 $p_i$ 和 $t_i$。试用 TabPFN 权重 $a$ 时，该记录的误差为
\begin{equation}
y_i-\big[a p_i+(1-a)t_i\big]=(y_i-t_i)-a(p_i-t_i).
\end{equation}
因此，融合是否减少一条记录的绝对误差，取决于分支预测差 $p_i-t_i$ 对树分支残差 $y_i-t_i$ 的修正。以 $I_g$ 表示内层折外记录中属于母体 $g$ 的行，按母体汇总该绝对误差，得到选权问题
\begin{equation}\label{eq:selectblend}
\hat a=\arg\min_{a\in\{0,\,0.25,\,0.5,\,0.75,\,1\}}
\frac{1}{G_{\rm in}}\sum_{g=1}^{G_{\rm in}}
\frac{1}{|I_g|}\sum_{i\in I_g}|(y_i-t_i)-a(p_i-t_i)|.
\end{equation}
离散网格仅需比较五项候选，两个端点对应单分支。选定配置和权重后，以整个外层训练集分别拟合两分支，在外层验证母体上输出
\begin{equation}\label{eq:newblend}
\hat y(x)=\hat a f_{\rm PFN}(x)+(1-\hat a)f_{\rm ET}(x).
\end{equation}
每折由种子 2026 选定配置和权重，再以种子 2026、17、42 分别拟合并评分。部署时在完整开发集上重复三折选参和选权，得到 23 维、2 成员的 TabPFN 与 869 维、$q=1$ 的 ExtraTrees，权重分别为 0.25 和 0.75。
